% ─── Cost-optimal pathway and country-specific heating mix ───────────────────
This section derives the 29-country least-cost residential heating pathway through
three steps: a source audit of the techno-economic inputs, a country-specific
reconstruction of the 2025 residential heating mix, and a least-cost
optimisation that \emph{derives} the per-country technology
mix from source-audited levelised costs instead of imposing it. The headline
result is that cost-optimal decarbonisation rests on district heat and heat pumps
at an approximately 53 to 45~per cent split, and that the depth of the scope-1
emissions cap barely moves the system cost, with 0.05~per cent separating the
$-75$, $-90$ and $-100$~per cent targets (Figure~\ref{fig:costopt_pathway}). The
cap is nonetheless not slack everywhere, and in the handful of countries where it
binds it does so at a substantial shadow price. Where this endogenous pathway is read
alongside the four exogenous-share scenarios, and what the comparison implies
for RQ1, is taken up in Section~\ref{sec:results:costoptcompare}; here the task is to
establish the least-cost pathway itself and to test its robustness.

\subsection{Source-Audited Techno-Economics}
Every capital-cost input is taken from three public catalogues: the Danish
Energy Agency technology data catalogue \citep{DEA2023technology}, the International
Energy Agency heat-pump cost work \citep{IEA2022heatpumps} and European Heat Pump
Association market data \citep{EHPA2024market}. Where they differ we adopt
the Danish Energy Agency central figure, because it is the most complete
and conservative of the three for the combustion and resistance technologies; the IEA
and EHPA sources are heat-pump-specific and corroborate the heat-pump inputs, not the
boiler inputs. Two of the resulting choices carry weight.

The first is the combustion boilers, which enter at Danish
Energy Agency central values (\euro{}420/kW gas, \euro{}450/kW oil, \euro{}950/kW
biomass and \euro{}200/kW resistance). The second is the hydrogen boiler, for which
none of the three catalogues reports
a residential input; we therefore adopt the Danish Energy Agency central
value of \euro{}600/kW, the most conservative public figure for this nascent technology.
Alongside the capital costs, the air-source heat-pump seasonal coefficient of
performance is the field value of 3.3 rather than the regulatory average-climate floor
of 3.0, and residential gas and electricity prices are taken from
Eurostat \citep{Eurostat2025prices}. Table~\ref{tab:costopt_audit} itemises the audit by
technology.

\begin{table}[htbp]
  \centering
  \caption{Source-audited technology capital costs.}
  \label{tab:costopt_audit}
  \footnotesize
  \begin{tabularx}{\textwidth}{@{}l c c c c X@{}}
    \toprule
    Technology & DEA 2023 & IEA 2022 & EHPA 2024 & Adopted & Basis \\
    \midrule
    Gas boiler        & \euro{}420/kW & n.r. & n.r. & \euro{}420/kW & DEA central; mature technology. \\
    Oil boiler        & \euro{}450/kW & n.r. & n.r. & \euro{}450/kW & DEA central; mature technology. \\
    Biomass boiler    & \euro{}950/kW & n.r. & n.r. & \euro{}950/kW & DEA central. \\
    Resistance        & \euro{}200/kW & n.r. & n.r. & \euro{}200/kW & DEA central; lowest-capital option. \\
    Hydrogen boiler   & \euro{}600/kW & n.r. & n.r. & \euro{}600/kW & No catalogue reports a residential input; DEA central, conservative for a nascent technology. \\
    Air-source HP     & corroborates  & corrob. & corrob. & SCOP 3.3 & EHPA/Eurovent field value, against a regulatory average-climate floor of 3.0. \\
    \bottomrule
  \end{tabularx}
  \\[2pt]{\footnotesize\textit{n.r. = not reported by that source. Source: Authors' contribution.}}
\end{table}

The same audit benchmarks the 2050 scenario technology shares against the EU
Reference Scenario 2024 \citep{EC2024RefScenario}, REPowerEU
\citep{REpowerEU2022}, Gas for Climate \citep{GasForClimate2050} and the
hydrogen-in-buildings consensus \citep{ICCT2023h2heating, H2ScienceCoalition},
and is documented in full in the supplementary material. One material reframing follows. The
high hydrogen share assumed under the H2 Push scenario is retained only as a
deliberate high-side stress test, since the mainstream literature places the
residential hydrogen share below 1 to 2~per cent.

\subsection{Country-Specific 2025 Heating Mix}
Each country's 2025 residential heating mix is estimated over the eight model
technologies from two independent bases and cross-checked. The first is an
\emph{energy} basis, the share of household final energy delivered by each fuel,
drawn from Eurostat disaggregated household final energy by fuel
\citep{Eurostat2024nrgdhhq}. The second is a \emph{dwelling} basis, the share of
dwellings equipped with each technology, drawn from national statistics,
Odyssee-Mure \citep{OdysseeMure} and the European Heat Pump Association. The two
bases carry different units, energy delivered against dwellings equipped, and so
bound the quantity the model actually needs from opposite sides.

The model allocates useful heat and not delivered energy, and the two bases
bracket that quantity. The energy basis over-weights inefficient technologies,
such as wood, oil and coal, which burn more delivered energy per unit of useful
heat; the dwelling basis over-weights technologies that deliver more useful heat
per unit of input, such as heat pumps and district heat. The per-country starting
mix is therefore taken as the mean of the two. As an illustrative example, a
country whose energy basis reads 65~per cent gas while its dwelling basis reads
45~per cent gas would enter the model at the 55~per cent mean; the true useful-heat
share lies between the two bounds by construction.

The divergence between the two bases at country level is retained as a validation
diagnostic and reported per country in the supplementary material; it is
widest where efficient and inefficient technologies coexist in similar shares and
narrowest in single-technology stocks. The resulting mixes are strongly
heterogeneous. The United Kingdom and the Netherlands begin at approximately 80~per cent
gas; Sweden and Denmark are district-heat dominated, at approximately
48~per cent and 66~per cent respectively; the Baltic states rely on a combination
of biomass and district heat; and Cyprus and Malta, which have no gas grid, heat
largely with reversible air-to-air units.

The starting mix is reconstructed per country rather than taken as a single EU
average, and this matters because the optimisation operates on each country's own
prices, feasibility limits and
baseline. A single EU-average mix would erase exactly the cross-country
heterogeneity, from gas-dominated to district-heat-dominated stocks, that
determines where each technology is least-cost. The per-country mix is used both
within the least-cost program and as the 2025 base mix from which the four multi-lever
scenarios (Section~\ref{sec:method}) interpolate.

\subsection{The Least-Cost Optimisation}
Whereas the four core scenarios impose exogenous 2050 technology shares,
the least-cost program solves for the cheapest mix. It is a linear program that
minimises the discounted, period-weighted system cost of the residential heating
pathway from 2025 to 2050, where $w_y$ is a period weight, $(1+r)^{-(y-2025)}$
discounts future years at rate $r$, $\mathrm{LCOH}_{c,t,y}$ is the levelised
cost of heat for technology $t$ in country $c$ and year $y$, $D_{c,y}$ is the
useful-heat demand and $x_{c,t,y}$ is the technology share. The objective is
\begin{equation}
\min \sum_{c,t,y} w_y\,(1+r)^{-(y-2025)}\,\mathrm{LCOH}_{c,t,y}\,D_{c,y}\,x_{c,t,y},
\label{eq:costopt_obj}
\end{equation}
minimised over technology shares $x_{c,t,y}$ for country $c$, technology $t$ and
milestone year $y$, subject to a set of constraints. A demand balance requires the shares
in each country and year to serve the full useful-heat demand,
\begin{equation}
\sum_{t} x_{c,t,y} = 1,\qquad x_{c,t,y}\ge 0,
\label{eq:costopt_demand}
\end{equation}
with a 2025 anchor fixing $x_{c,t,2025}$ to the country's observed mix. Feasibility
ceilings cap the deployable share of district heat, ground-source heat pumps,
hydrogen and sustainable biomass,
\begin{equation}
x_{c,t,y} \le \bar{x}_{c,t},\qquad t\in\{\text{DH},\text{GSHP},\text{H}_2,\text{biomass}\},
\label{eq:costopt_ceiling}
\end{equation}
and a stock-turnover limit bounds the year-on-year change in each share by the
fraction of the boiler stock that reaches end of life within the milestone interval,
\begin{equation}
\lvert x_{c,t,y}-x_{c,t,y-1}\rvert \le \delta_{c,t},
\label{eq:costopt_turnover}
\end{equation}
alongside the boiler-ban phase-out drawn from the policy layer. Finally, a
per-country emissions cap limits 2050 scope-1 emissions,
\begin{equation}
\sum_{t} e_{t}\,D_{c,y}\,x_{c,t,y}
\le \left(1-\kappa\,\frac{y-2025}{25}\right) E_{c,2025}^{\text{scope-1}},
\label{eq:costopt_cap}
\end{equation}
where $e_t$ is the on-site emission factor of technology $t$,
$E_{c,2025}^{\text{scope-1}}$ the country's 2025 scope-1 baseline, and $\kappa\in[0,1]$
the required reduction fraction relative to that baseline at 2050. The cap glides
linearly from the country's own 2025 baseline to $(1-\kappa)$ of it at 2050 rather than
binding at the terminal value in every year; imposing the terminal cap throughout is
infeasible at $\kappa=0.90$ given the turnover limit of
Equation~\eqref{eq:costopt_turnover}. The cap is applied to scope-1
emissions only, that is, to on-site gas and oil combustion, consistent with the
zero-emission building definition of the recast Energy Performance of Buildings
Directive \citep{EPBD2024}, so that the exogenous decarbonisation of the grid and of
district heat is not double-counted. The scope-1 boundary ensures the model does not
double-count exogenous grid decarbonisation or district-heat clean-up. It reflects the
fact that building-level policy levers, namely boiler choice and envelope depth,
directly control on-site emissions, and grid decarbonisation is an EU-wide system outcome that the building sector does not set. Three ambition variants set $\kappa$ to
$0.75$, $0.90$ and $1.00$, corresponding to scope-1 cuts of $-75$~per cent, $-90$~per
cent and $-100$~per cent against each country's own 2025 baseline. The objective uses
the source-audited levelised cost of heat, with the ETS2 carbon price embedded, so that
the cost-optimal allocation reflects both the corrected techno-economics and carbon
pricing.

The shadow price of each country's cap is the implied carbon price that the country
would require in order to meet its target.\footnote{Four countries, Estonia, Malta,
Portugal and Sweden, return a degenerate dual because their absolute scope-1 baseline is
small enough that a tiny tonnage divides into a large per-tonne value that carries no
meaning as a marginal abatement cost. Their shadow prices are undefined and they are
excluded from the binding-country count throughout. The largest well-defined value
anywhere on the path is Cyprus at about \euro{}207/tCO$_2$ in 2030.} To
isolate the least-cost technology mix independent of assumptions about envelope retrofit
depth, the demand $D_{c,y}$ is fixed at the Stated Policies LMDI trajectory
(Section~\ref{sec:method}), using the central per-country envelope rate.

This choice reflects the modelling premise. The least-cost program answers which heating
technology is cost-optimal at a given demand level, so the demand path is held at a
moderate, policy-consistent trajectory and not the deeper Net Zero path, whose own
demand reduction is itself an optimisation outcome that would prejudge the
supply-versus-demand margin the optimisation is meant to test. The robustness test
below, which co-optimises retrofit depth, confirms that supply-side switching is the
least-cost margin and that this choice does not drive the result. Accordingly the 2050
EU demand of approximately 3,480~TWh tracks Stated Policies.

The 29-country least-cost 2050 mix is approximately 53~per cent district heat and
45~per cent heat pump, with hydrogen at 0.34~per cent and residual gas at
1.4~per cent. The shares are expressed on a useful-heat basis, that is, as a
fraction of the building's physical heat demand, the convention used for every
share in this paper. The ordering of the two leading technologies
follows directly from charging equipment capital against peak sizing, not against annual output, which raises the heat pump's levelised cost by more than it
raises any boiler's, since a heat pump is sized to the same peak whether or not it
is efficient.

The composition of the mix is now governed almost entirely by the deployment
ceilings and not by the cost ranking, and the two have come apart. On 2050
levelised cost the cheapest single technology is district heat in 25 countries, the
hydrogen boiler in 2 and a heat pump in 2. Heat pumps nonetheless
take 45~per cent of EU useful heat, because air-source units are the only major
option in the program that carries no feasibility ceiling of its own
(Eq.~\eqref{eq:costopt_ceiling}), so they absorb whatever demand the ceiling-limited
technologies cannot. The corollary matters for the paper's central question.
Hydrogen's 0.34~per cent share is a ceiling-and-turnover outcome, not a cost
outcome, and it is not evidence that the least-cost program
rejects hydrogen on price. The program does not reject hydrogen on price; the
feasibility ceiling prevents it from selecting more. That ceiling should be read for
what it is, and it is stronger than a low cap. It is 3~per cent in the Netherlands and
2~per cent in the United Kingdom and \emph{zero in the other 27 countries}, so in 27
markets the optimisation is not permitted to select any hydrogen at all. Denmark is the
sharpest case, because the hydrogen boiler is one of the two markets where it is the
cheapest single 2050 technology and its ceiling is zero. The ceiling encodes the
mainstream literature's sub-2~per cent residential share and the physical pace at which
a distribution network could be converted, which is defensible, but a near-zero
hydrogen share in this optimisation is an assumption and not a finding.

Within the heat-pump category the optimisation now favours air-source units almost
exclusively, at 44.4~per cent of useful heat against 1.05~per cent ground-source, so
the EU mix moves from a fossil-dominated 2025 baseline to a 2050 mix led jointly by
district heat and air-source heat pumps (Figure~\ref{fig:costopt_pathway}).
Ground-source units sit well inside their feasibility ceiling; at 1.05~per cent they sit well below it, because their higher capital cost is charged in full against peak sizing. District heat
takes the balance in dense stock and in areas served by an existing network.

The second result of the optimisation concerns the price of carbon. The cost of \emph{ambition} is
negligible. The discounted present-value system cost is
\euro{}7{,}426.3~bn~EUR-yr at the $-75$~per cent cap, \euro{}7{,}428.4~bn at
$-90$~per cent and \euro{}7{,}429.8~bn at $-100$~per cent, a spread of
0.05~per cent across the three depths, so tightening the scope-1 target all the way
to net zero on-site combustion costs almost nothing at the audited costs.
Co-optimising retrofit alongside the supply mix lowers the $-90$~per cent present value by \euro{}24.8~bn, and the $-100$~per cent one by \euro{}24.7~bn, 0.33~per cent in each case. That is small against the level and about seven times the 0.05~per cent the cap depth moves, so supply-side switching remains the cheap abatement margin. The renovation the least-cost solution does take, it takes at the interim milestones and lets lapse by 2050, at 2.8~per cent of EU useful heat in 2030 and 5.8~per cent in 2040 (Appendix~\ref{sec:costopt}).

The \emph{shadow price}, however, is not uniformly zero, and where it is not the policy
reading changes. The cap binds in 4 of the 29 countries in 2050, namely
Cyprus, Croatia, Poland and the United Kingdom, at an implied carbon price of up to
\euro{}81/tCO$_2$, and it binds in seven countries at some point along the 2025 to
2050 path, the three additions being Spain, Hungary and Romania. The four
degenerate cases noted in the optimisation set-up above are excluded from both
counts. A
cap that binds in a handful of markets is a materially different policy statement
from a cap that is slack everywhere. The first says that ETS2 pricing carries most
of the EU to its target but leaves an identifiable minority needing a carbon price
several times higher, or an equivalent non-price instrument, to close the last
increment. A cap that were slack everywhere would instead say that no supplementary
instrument is needed anywhere, and that is not what the optimisation gives.

The implication of the cost-optimal result is therefore narrow, and should be stated
narrowly. For most of the EU, residential heating decarbonisation
is cost-effective at current ETS2 pricing, in the sense that the cap does not bind
and no carbon price supplement is required in the model to reach the target. For a
minority of countries it is not, and the binding shadow price there is large enough
that the policy question in those markets is a live one. This finding rests on the
two pillars documented above, the audited techno-economics and the embedding of
ETS2 in the levelised-cost ranking, and it is conditional throughout on the
optimisation's perfect-foresight least-cost framing.

\paragraph{Robustness of the cost-optimal mix to cap ambition.} The perturbation
that matters most for the policy reading is the depth of the emissions cap, and
Table~\ref{tab:costopt_robust} reports it directly. The present-value system cost
moves by 0.05~per cent across the three depths, and district heat and hydrogen hold
their 2050 shares across them while the adjustment falls entirely on the heat-pump
against residual-gas margin, which runs 43.3 to 45.4 to 46.8~per cent against 3.5 to 1.4
to zero. The \emph{cost} of ambition is negligible in that sense.
What is not stable is the binding status. The cap is slack in most
countries at every depth, but it binds in 4 countries in 2050 and in 7 at some
point on the path. Where it binds in 2050 the implied carbon price reaches
\euro{}81/tCO$_2$, and it is not monotone along the path. Six countries bind in 2030
at up to \euro{}207/tCO$_2$ in Cyprus, three in 2040 at up to \euro{}37, and four in
2050. The shadow-price column is not a uniform zero, and the
near-term figures are the larger ones.

\begin{table}[htbp]
  \centering
  \caption{The cost of ambition and the cap's binding status across the scope-1 caps.}
  \label{tab:costopt_robust}
  \small
  \begin{tabularx}{\textwidth}{@{}l c c X@{}}
    \toprule
    Cap ambition & Present-value system cost & Spread & Binding status \\
                 & (bn EUR-yr)               &        &                \\
    \midrule
    $-75$\,\% scope-1  & 7{,}426.3 & reference        & Binds in 4 of 29 at up to \euro{}81/tCO$_2$. \\
    $-90$\,\% scope-1  & 7{,}428.4 & $+0.03$\,\%      & Binds in the same 4 at the same shadow price. \\
    $-100$\,\% scope-1 & 7{,}429.8 & $+0.05$\,\%      & Binds in the same 4 even at net zero on-site combustion. \\
    \bottomrule
  \end{tabularx}
  \smallskip
  \raggedright\small
  The spread column is measured against the $-75$~per cent cap. In the headline
  program the cap binds in 4 of the 29 countries in 2050 (Cyprus, Croatia, Poland
  and the United Kingdom) at an implied carbon price of up to \euro{}81/tCO$_2$,
  and in seven countries at some point on the 2025 to 2050 path (adding Spain, Hungary
  and Romania). The four degenerate cases (Estonia, Malta, Portugal and Sweden) are
  excluded throughout.
  \\[2pt]{\footnotesize\textit{Source: Authors' contribution.}}
\end{table}

Two further perturbations bear on the mix alone. The first is the
feasibility ceiling on ground-source heat pumps, which does not bind. With
equipment capital charged against peak sizing, ground-source
units carry their higher capital in full, so the program selects only
1.05~per cent of useful heat from them, well inside any ceiling in the tested range.
The ground-source ceiling is accordingly not a binding assumption of the
result, and the per-ceiling sweep is documented in the supplementary material. The
biomass ceiling is likewise not binding, since the linear program selects
approximately 0~per cent biomass across a ceiling range of 0.10 to 0.30. The
second is the accounting boundary of the cap, scope-1 against a scope-2-inclusive
total, which is reported in the supplementary material and discussed as a limitation in
the emissions-boundary subsection of Appendix~\ref{sec:limitations}.

Finally, as a robustness test outside the central parameter set, we relax the
fixed-demand assumption itself by endogenising renovation depth. A linear negawatt
retrofit option avoids useful heat at a country-specific cost drawn from the medium-depth band of the European Commission's renovation-activity study \citep{Esser2019renovation}, set at \euro{}154/m$^2$ for an approximately 41~per cent saving and annualised over a 30-year envelope life, which gives a cost of \euro{}117 to
1{,}187/MWh of useful heat avoided, cheapest in cold, high-intensity stock.
Co-optimising retrofit with the supply mix leaves the 2050 mix essentially unchanged, and the least-cost solution selects 0.0~per cent retrofit in 2050 even at the net-zero cap. That endpoint is not the whole story, and reporting it alone would misstate the program's behaviour. Retrofit is selected at the interim milestones and abandoned by 2050: 2.8~per cent of EU useful heat in 2030, in six countries, and 5.8~per cent in 2040, all of Germany's. Every one of those shares sits exactly on its ramp limit, so the program wants more retrofit at those dates than the 2~percentage-point-a-year build rate allows. The ramp is one-sided, with no floor holding the stock in place, so the program builds envelope for the interim caps and lets it lapse when clean supply has caught up. The \euro{}24.8~bn present-value improvement the variant delivers is the value of that interim renovation, and not a residual.

The cheapest negawatt costs about \euro{}117/MWh of useful heat avoided, and the
heat-pump levelised cost spans about \euro{}65 to 197/MWh across the 29 countries,
so the two ranges overlap and the inequality does not hold universally. In the dearest
heat-pump markets a deep retrofit is cheaper per avoided megawatt-hour than the supply
switch it would displace, and the zero-retrofit result survives on the median alone. This still confirms the modelling choice to fix the least-cost program's
demand at the Stated Policies trajectory, because the optimum is unchanged, but it
should not be presented as a strict cost dominance. It is consistent in that
qualified form with the Joint Research Centre finding \citep{Filippidou2019} that deep
envelope renovation is cost-optimal only in specific contexts, namely central-European,
high-demand stock, not as a general abatement lever.

\subsection{Limitations of the Cost-Optimal Model}
The robustness tests above show that the headline mix and the negligible cost of
ambition survive the central parameter choices, while the zero implied carbon price
does not. Three structural assumptions of the linear program further qualify how
its results, and in particular the zero-retrofit outcome, should be read. They are
stated here in full so that the cost-optimality headline is interpreted correctly;
none overturns the mix, but each bounds its reach.

First, the baseline program assumes \emph{perfect foresight}. It minimises a
discounted system cost computed against known future levelised costs, so agents in
effect know the entire 2025 to 2050 price and carbon trajectory in advance, and
under that assumption the cost-optimal solution selects approximately zero per cent
envelope retrofit because, at the audited costs, renovation is dearer per
megawatt-hour than switching supply in most of the stock. The natural
concern is that the near-zero retrofit is an artefact of this foresight. It is not.
Relaxing perfect foresight to a \emph{myopic rolling-horizon} optimisation, in which
each milestone year is solved on the costs visible at the time with a one- or
two-period look-ahead, not the full trajectory, leaves the result unchanged. The 2050 retrofit share stays at 0.0~per cent, and the heat-pump, district-heat and hydrogen mix is identical in all 29 countries across
the foresight modes at each target, reproducing the headline mix reported above with hydrogen at 0.34~per cent. Under the one- and two-period look-ahead cases the realised present-value cost matches the perfect-foresight optimum of about \euro{}7{,}430~bn.

The zero-look-ahead snapshot, which re-optimises each
milestone year in isolation with no inter-period stock memory, is the exception. It
\emph{diverges} on total cost, because the isolated single-year program does not
carry the cumulative replacement obligation, yet still returns the same 0.0~per cent 2050 retrofit and the identical technology mix, so the divergence is in the cost accounting
only and the mix verdict is intact. The result is foresight-invariant because the pathway is
bound by stock turnover and the feasibility ceilings, with no room for intertemporal arbitrage to bind.

Retrofit is not, however, cost-dominated at every horizon. The cheapest negawatt retrofit costs about \euro{}117/MWh of useful
heat avoided, against a heat-pump levelised cost of heat spanning about
\euro{}65 to 197/MWh, so the two ranges overlap and retrofit is cheaper than a heat
pump in the dearest heat-pump markets. The solved program still selects 0.0~per cent retrofit in 2050, and it does so because district heat, at a 2050 EU-median levelised cost of \euro{}96.5/MWh, has by then displaced the cheapest negawatts, while the ramp limit is what holds the interim shares down. Renovation is not dearer than every supply option everywhere, and the program takes it where the interim caps force early abatement. The zero-retrofit outcome therefore survives as a result of the optimisation. The caveat on it
is cost-and-constraint-conditioned rather than perfect-foresight-conditioned, and it
carries a country caveat as well. Real building owners
still face genuine uncertainty over future energy and carbon prices, which can
justify precautionary envelope investment that neither optimiser values, so the
modelled retrofit share remains a lower bound on real-world behaviour.

Second, the model charges \emph{no switching cost beyond annualised capital}. It
applies only the stock-turnover constraint of
Equation~\eqref{eq:costopt_turnover}, which limits the pace of replacement to the
rate at which boilers reach end of life but does not penalise replacing equipment
before its useful life is over. The economic cost of premature scrappage, and any
transaction or disruption cost of switching technology, is not represented, so the
modelled pace of transition is an upper bound on what an equivalent
switching-cost-aware optimisation would deliver.

Third, the model holds the \emph{demand path fixed} at the Stated Policies LMDI
trajectory instead of co-optimising it with supply. This is a deliberate choice,
made so that the least-cost program answers which heating technology is cost-optimal at
a given demand level and not how much to invest in renovation; the
endogenous-retrofit variant above confirms that relaxing it does not change the
mix. It nonetheless means the reported pathway is the least-cost way to serve a fixed demand, and the deeper demand
reductions of the Net Zero path are excluded by construction. These limitations do
not invalidate the cost-optimal result, which is robust across the sweeps reported
above, but they delimit it. The result is a least-cost \emph{supply} pathway under
foresight, and is best read as one bound on real-world behaviour.

\begin{figure}[htbp]
  \centering
  \caption{The EU least-cost heating pathway under the $-90$~per cent scope-1 cap. Each
    band is a technology's share of EU useful heat, so the vertical axis runs from 0 to 1
    and the bands sum to it in every year. Hydrogen carries a band of 0.3~per cent, which
    is not visible at this scale, and that absence is the figure's point.}
  \includegraphics[width=0.8\textwidth]{paper/P17_costopt_pathway.pdf}
  \label{fig:costopt_pathway}
  \\[2pt]{\footnotesize\textit{Source: Authors' contribution.}}
\end{figure}

Country-specific outcomes under the $-90$~per cent cap are not uniform, and the
non-uniformity is where the policy content sits. With ETS2 embedded in levelised
costs, most countries reach the target with the cap slack and an implied carbon
price of zero, but 4 of the 29 reach 2050 with the cap binding, at up to
\euro{}81/tCO$_2$, and 7 have a binding cap at some point on the path. The
cost-optimal district-heat and heat-pump shares vary with feasibility and climate
around the EU mix of Figure~\ref{fig:costopt_pathway}. The per-country shares and
shadow prices are tabulated in the supplementary material.
